\section{Eksperymenty i wyniki}
\label{sec:eksperymenty}
Podrozdział przedstawia pomiary wykonane za pomocą opisanego wyżej systemu. Porównywane warianty są
konfiguracjami tego samego grafu, przedstawionego w podrozdziale \ref{sec:sciezka}. Przedmiotem
porównania są dwie architektury przetwarzania zapytania głosowego: klasyczna kaskada, w której
nagranie jest najpierw transkrybowane, a dopiero powstały tekst embeddowany w przestrzeni
wektorowej, oraz proponowane rozwiązanie, w którym reprezentacja latentna modelu ASR jest
dopasowywana wprost do przestrzeni modelu embeddingowego. Obu towarzyszy trzecia konfiguracja,
pozbawiona ścieżki głosowej i pełniąca rolę punktu odniesienia, a nie rozwiązania poddanego ocenie.
Pomiary prowadzone są na dwóch zestawach danych opisanych w podrozdziale~\ref{sec:zbiory}.
Porównywane konfiguracje przetwarzają w obrębie jednego zestawu ten sam zbiór pytań, więc różnice
liczone są parami.

\subsection{Metodologia}

\subsubsection{Zbiory danych}
\label{sec:zbiory}
Podstawą porównania jest zbiór PubMedQA \cite{jin2019pubmedqa}, w którym każde pytanie ma przypisany
artykuł źródłowy, pełniący rolę jedynego dokumentu istotnego. Zbiór występuje w dwóch
konfiguracjach: anotowanej przez ekspertów, liczącej tysiąc pytań, oraz etykietowanej automatycznie,
obejmującej ponad dwieście tysięcy pytań. Wybrano konfigurację anotowaną, mimo że jest o dwa rzędy
wielkości mniejsza, ponieważ w konfiguracji automatycznej ponad dziewięćdziesiąt procent odpowiedzi
to ,,tak'', a klasa ,,może'' praktycznie nie występuje. Trafność decyzji mierzona na takim materiale
nie odróżniałaby systemu od odpowiadania stałą wartością.

Eksperymenty prowadzi się na dwóch zestawach, z których każdy odpowiada na inne pytanie. Zestaw
główny to całość konfiguracji anotowanej: tysiąc pytań i tysiąc dokumentów, w całości
anglojęzycznych. Służy on porównaniu samych architektur przetwarzania, ponieważ daje moc testu oraz
dziesięciokrotnie większy korpus, w którym trafienie w dokument istotny jest odpowiednio
trudniejsze. Zestaw językowy liczy sto pytań, jest ścisłym podzbiorem poprzedniego i występuje w
trzech wariantach językowych; służy zmierzeniu kosztu zadania pytania po polsku oraz ocenie
odpowiedzi generowanych z pełnych artykułów.

Podział ten wynika z natury badanego zjawiska. Różnica między kaskadą a dopasowaniem reprezentacji
latentnych bierze się z tego, że pierwsza przepuszcza sygnał przez pośrednią postać dyskretną, a
druga nie; mechanizm ten nie zależy od języka zapytania. Rozstrzygnięcie, która architektura jest
lepsza, wymaga zatem przede wszystkim mocy statystycznej, którą daje zestaw większy. Osobnym
pytaniem jest to, ile jakości ubywa, gdy zapytanie pada po polsku, i na nie odpowiada zestaw
mniejszy.

Zestaw językowy liczy sto pytań, a nie tysiąc, z dwóch powodów naraz. Po pierwsze wariant z korpusem
polskim wymaga przetłumaczenia każdego dokumentu, a tłumaczenie tysiąca abstraktów przekraczało
dostępne zasoby. Po drugie ocena generacji wymaga pełnych tekstów artykułów, te zaś są dostępne
wyłącznie dla podzbioru: spośród tysiąca artykułów konfiguracji anotowanej pełny tekst udaje się
pobrać z otwartego repozytorium PMC dla stu piętnastu, czyli dla 11,5 procent. Wybrano sto pytań
spełniających oba warunki naraz, dzięki czemu ten sam zestaw obsługuje część językową i część
generacyjną.

Korpus zestawu językowego przygotowano w trzech wariantach, aby rozdzielić dwa efekty, które w
pojedynczym pomiarze pozostawałyby ze sobą splecione. Wariant pierwszy zestawia pytania polskie z
korpusem angielskim i odpowiada sytuacji realnej, ponieważ literatura biomedyczna jest w
przeważającej mierze anglojęzyczna. Wariant drugi zestawia pytania polskie z korpusem
przetłumaczonym i pozwala zmierzyć, ile jakości traci samo przejście między językami. Wariant
trzeci, w całości anglojęzyczny, pokazuje koszt użycia języka polskiego niezależnie od modalności.
Podział ten odpowiada wprost ograniczeniom omówionym w podrozdziale \ref{sec:motywacja}. Wszystkie
trzy warianty obejmują te same sto pytań i te same sto dokumentów, rozpoznawane po wspólnych
identyfikatorach, więc różnice liczone są parami, pytanie po pytaniu.

Tłumaczenie na język polski wykonano modelem językowym, osobno dla treści pytania, tytułu dokumentu,
abstraktu i odpowiedzi referencyjnej; etykiety decyzji pozostawiono w postaci oryginalnej, ponieważ
ocena porównuje je z literałami zbioru. Pełne teksty artykułów pozostają wyłącznie anglojęzyczne,
zatem ocena odpowiedzi generowanych z pobranego kontekstu przebiega zawsze na materiale angielskim,
także wtedy, gdy pytanie zadano po polsku. Charakterystykę obu zbiorów zestawiono w
tabeli~\ref{tab:zbior}.

\begin{table}[h]
\centering
\small
\begin{tabular}{@{}L{6.4cm}L{2.5cm}L{2.5cm}L{2.5cm}@{}}
\toprule
& \textbf{Zestaw główny} & \multicolumn{2}{c}{\textbf{Zestaw językowy}} \\
\cmidrule(lr){2-2}\cmidrule(lr){3-4}
\textbf{Wielkość} & angielski & angielski & polski \\
\midrule
Liczba pytań                              & 1000          & 100        & 100 \\
Liczba dokumentów w korpusie              & 1000          & 100        & 100 \\
Dokumenty istotne na pytanie              & 1             & 1          & 1 \\
Rozkład odpowiedzi (tak / nie / może)     & 552/338/110   & 56/36/8    & 56/36/8 \\
Długość pytania (mediana, słowa)          & 13            & 12,5       & 13 \\
Długość abstraktu (mediana, słowa)        & 201           & 209        & 206 \\
Długość pełnego artykułu (mediana, słowa) & ---           & 3286       & --- \\
\bottomrule
\end{tabular}
\caption{Charakterystyka zbiorów użytych w eksperymentach. Zestaw językowy jest ścisłym podzbiorem
głównego, a jego wersje językowe obejmują te same pytania i dokumenty, rozpoznawane po wspólnych
identyfikatorach. Pełne teksty artykułów są dostępne wyłącznie w wersji angielskiej zestawu
językowego.}
\label{tab:zbior}
\end{table}

Oba korpusy są małe wobec zastosowania rzeczywistego, w którym system przeszukuje zbiór o kilka
rzędów wielkości większy, więc bezwzględne wartości metryk rankingowych należy czytać jako
optymistyczne. Dotyczy to zwłaszcza korpusu stu dokumentów, gdzie nawet losowy ranking trafia w
dokument istotny co dziesiąte pytanie przy progu odcięcia równym dziesięć. Nie unieważnia to
porównania, ponieważ wszystkie konfiguracje pracują na tym samym korpusie, a przedmiotem oceny jest
różnica między nimi, nie poziom bezwzględny. Ponieważ zestaw językowy jest podzbiorem głównego,
wynik konfiguracji anglojęzycznej można policzyć na obu i sprawdzić, czy setka zachowuje się jak
tysiąc; rozbieżność wskazywałaby, że efekt przypisany językowi bierze się z doboru podzbioru.

Ścieżka audio nie pochodzi z nagrań, lecz jest syntezowana w czasie wykonania przez węzeł syntezy
mowy. Wybór ten zapewnia, że wszystkie porównywane konfiguracje pracują na identycznym sygnale
wejściowym, więc różnice w wynikach nie mogą pochodzić z różnic między nagraniami. Ceną jest
ograniczona zmienność mowy: sygnał syntetyczny nie obejmuje pełnego zakresu cech osobniczych mówcy,
wad wymowy ani zakłóceń otoczenia. Wnioski o odporności należy zatem odnosić do zmienności
wprowadzonej celowo, opisanej w podrozdziale \ref{sec:odpornosc}, a nie do zmienności naturalnej.

\subsubsection{Porównywane architektury}
Porównaniu podlegają trzy konfiguracje. Pierwsza pomija ścieżkę głosową i kieruje do wyszukiwania
tekst pytania. Nie jest ona rozwiązaniem poddanym ocenie, lecz wyznacza górną granicę jakości
osiągalnej przy danym korpusie i danym modelu embeddingowym. Bez tego odniesienia nie sposób
rozstrzygnąć, czy różnica między pozostałymi konfiguracjami wynika z architektury, czy z trudności
samego zadania wyszukiwania. Druga konfiguracja to kaskada, w której nagranie przechodzi przez model
ASR, a powstała transkrypcja przez model embeddingowy. Trzecia to dopasowanie reprezentacji
latentnych, w którym wyjście enkodera ASR rzutowane jest wprost do przestrzeni modelu embeddingowego
przez adapter wytrenowany w rozdziale \ref{sec:latentTranslation}.

\todo{Uzupełnić listę wykorzystanych modeli ASR i embeddingowych wraz z uzasadnieniem doboru, opartym na wnioskach z rozdziałów \ref{sec:latent} i \ref{sec:latentTranslation}}

\subsubsection{Metryki}
Jakość wyszukiwania oceniana jest metrykami rankingowymi zdefiniowanymi w podrozdziale
\ref{sec:metryki}. Dobór metryk ogranicza jednak struktura zbioru: każde pytanie ma dokładnie
jeden dokument istotny, czyli artykuł, z którego powstało. Dla pojedynczego pytania skuteczność
wyszukiwania przyjmuje wtedy wartość zero albo jeden, precyzja na progu $k$ jest jej stałą
wielokrotnością, a średnia precyzja pokrywa się z odwrotnością rangi. Raportuje się zatem Recall@k
oraz NDCG@k dla progów $k \in \{1,5,10\}$, a także MRR liczoną dla całej listy wyników; Precision@k
i MAP pominięto, ponieważ na tych danych nie wnoszą nic ponad to, co zawierają pozostałe miary.
Metryki liczone są osobno dla każdego pytania, co jest warunkiem zastosowania testu sparowanego oraz
analizy przedstawionej w podrozdziale \ref{sec:prog}.

Jakość transkrypcji mierzona jest wskaźnikiem błędu na poziomie słów i znaków. Wielkości te dotyczą
wyłącznie konfiguracji kaskadowej, ponieważ pozostałe dwie transkrypcji nie wytwarzają. Nie są więc
przedmiotem porównania między architekturami, lecz zmienną objaśniającą w analizie zależności między
jakością transkrypcji a jakością wyszukiwania.

\subsubsection{Aparat statystyczny}
Każdą wielkość zbiorczą podaje się wraz z przedziałem ufności wyznaczonym metodą losowania ze
zwracaniem (ang.\ \textit{bootstrap}), a różnice między konfiguracjami ocenia się testem sparowanym,
ponieważ wszystkie konfiguracje przetwarzają ten sam zbiór pytań. Zmienność między pytaniami jest
znacznie większa niż różnice między architekturami, więc test dla prób niezależnych miałby na tych
danych zbyt małą moc, by wykryć badane różnice.

Porównanie architektur, analizę odporności, wyznaczenie progu jakości transkrypcji oraz pomiary
zasobów prowadzi się na zestawie tysiąca pytań, ponieważ wszystkie wymagają mocy testu. Efekt języka
oraz ocenę odpowiedzi generowanych z kontekstu prowadzi się na zestawie stu pytań, ponieważ tylko on
istnieje w wariantach językowych i tylko dla niego dostępne są pełne teksty. Każde porównanie opiera
się na pełnym komplecie par danego zestawu; sto par jest liczbą skromną, lecz przekracza próg,
poniżej którego framework oznacza wynik jako obarczony małą mocą.

\subsection{Porównanie architektur w warunkach czystych}
\label{sec:czyste}
Podrozdział przedstawia wyniki wyszukiwania dla nagrań niezakłóconych na zestawie tysiąca pytań.
Jest to zasadnicze porównanie niniejszej pracy: kaskada zestawiona z dopasowaniem reprezentacji
latentnych, obie odniesione do górnej granicy wyznaczonej przez konfigurację tekstową. Odniesienie
do tej granicy rozdziela dwa źródła utraty jakości, przejście na modalność mowy oraz wybór
architektury przetwarzania, których łączny efekt w pojedynczym pomiarze pozostawałby
nierozróżnialny.

\todo{Uzupełnić tabelę wyników oraz ich omówienie po wykonaniu przebiegów}

\subsection{Koszt zadania pytania po polsku}
\label{sec:jezyk}
Podrozdział powtarza porównanie architektur na zestawie stu pytań, w trzech wariantach językowych
opisanych w podrozdziale~\ref{sec:zbiory}, i na tej podstawie rozdziela utratę jakości wnoszoną
przez sam język od utraty wnoszonej przez modalność. Wariant w całości anglojęzyczny odtwarza
warunki podrozdziału~\ref{sec:czyste} na mniejszej próbie, więc pełni podwójną rolę: wyznacza punkt
odniesienia dla dwóch pozostałych wariantów, a zarazem pozwala sprawdzić, czy setka zachowuje się
jak tysiąc, z którego została wybrana.

Rozstrzygnięcia wymaga zwłaszcza różnica między wariantem jednojęzykowym polskim a międzyjęzykowym,
ponieważ wskazuje ona, jaka część utraty jakości bierze się z przekraczania granicy języka przez
samo zapytanie, a jaka z tego, że polszczyzna jest dla użytych modeli językiem gorzej
reprezentowanym niż angielski.

\todo{Uzupełnić zestawienie wyników dla trzech wariantów językowych oraz omówienie po wykonaniu przebiegów}

\subsection{Jakość odpowiedzi generowanych z pobranego kontekstu}
\label{sec:generacja}
Wyszukiwanie jest w zastosowaniu klinicznym środkiem, a nie celem: użytkownik oczekuje odpowiedzi,
nie listy dokumentów. Podrozdział sprawdza zatem, czy przewaga jednej architektury nad drugą,
zmierzona metrykami rankingowymi, przenosi się na jakość odpowiedzi generowanej z pobranego
kontekstu. Przejście to nie jest oczywiste, ponieważ model językowy potrafi odpowiedzieć poprawnie
mimo niedoskonałego rankingu, o ile potrzebny fragment znalazł się w kontekście, i odwrotnie ---
potrafi zawieść przy rankingu poprawnym.

Kontekstem są pełne artykuły, a nie abstrakty, ponieważ abstrakt rzadko zawiera dane liczbowe
potrzebne do uzasadnienia odpowiedzi. Jest to jednak kontekst wyłącznie angielski, więc konfiguracje
z pytaniem polskim pracują w układzie międzyjęzykowym również na etapie generacji. Mediana długości
artykułu przekracza trzy tysiące słów, czyli około siedemnastokrotność abstraktu, co przy kilku
pobranych dokumentach wyczerpuje typowe okno kontekstowe modelu językowego. Liczbę dokumentów
podawanych modelowi dobrano więc z uwzględnieniem tego ograniczenia.

Ocena obejmuje metryki opisane w podrozdziale \ref{sec:rag}: trafność decyzji wobec etykiety zbioru,
podobieństwo ROUGE do odpowiedzi referencyjnej, wskaźnik halucynacji, pokrycie kontekstu oraz oceny
sędziego automatycznego. Pokrycie kontekstu pełni tu rolę rozstrzygającą, ponieważ oddziela
niepowodzenie wyszukiwania od niepowodzenia generacji: jego niska wartość oznacza, że potrzebnej
treści w pobranych dokumentach nie było, więc winą za błędną odpowiedź nie można obciążać modelu
językowego.

\todo{Uzupełnić zestawienie metryk odpowiedzi oraz omówienie po wykonaniu przebiegów}

\subsection{Odporność na degradację sygnału}
\label{sec:odpornosc}
Podrozdział przedstawia jakość wyszukiwania jako funkcję stosunku sygnału do szumu, osobną krzywą
dla każdej z porównywanych architektur, a ponadto wpływ zmiany tempa mowy. Celem
jest sprawdzenie hipotezy o odmiennym charakterze degradacji obu rozwiązań. W kaskadzie błąd
transkrypcji usuwa termin z zapytania bezpowrotnie, ponieważ dalsze etapy nie mają już dostępu do
sygnału, z którego termin powstał. W dopasowaniu reprezentacji latentnych sygnał nie przechodzi
przez pośrednią postać dyskretną, więc oczekuje się degradacji łagodniejszej, rozłożonej na
wszystkie wymiary reprezentacji zamiast skupionej w pojedynczym utraconym słowie.

\todo{Uzupełnić wykresy jakości w funkcji stosunku sygnału do szumu oraz omówienie po wykonaniu przebiegów}

\subsection{Zależność jakości wyszukiwania od jakości transkrypcji}
\label{sec:prog}
Podrozdział zestawia jakość wyszukiwania z jakością transkrypcji, pytanie po pytaniu, i na tej
podstawie wyznacza próg jakości transkrypcji, poniżej którego kaskada przestaje przewyższać
rozwiązanie proponowane. Ponieważ konfiguracja latentna transkrypcji nie wytwarza, jej wyniki
przypisuje się wartości wskaźnika błędu osiągniętej przez kaskadę na tym samym pytaniu. Zestawienie
takie jest możliwe wyłącznie dzięki temu, że obie konfiguracje przetwarzają ten sam zbiór, a ich
wyniki pozostają powiązane identyfikatorem pytania.

Jest to wynik, który najbardziej bezpośrednio przekłada się na regułę doboru architektury. Znając
spodziewaną jakość transkrypcji w danym zastosowaniu, można z niego odczytać, która ścieżka
przetwarzania jest korzystniejsza, bez powtarzania całego porównania.

\todo{Uzupełnić wykres zależności oraz wyznaczoną wartość progu wraz z przedziałem ufności}

\subsection{Wykorzystanie zasobów i prędkość wyszukiwania}
\label{sec:zasoby}
Podrozdział przedstawia opóźnienie pełnej ścieżki przetwarzania, szczytowe zużycie pamięci karty
graficznej, przepustowość oraz rozmiar indeksu. Oczekiwana przewaga konfiguracji latentnej wynika z
pominięcia dekodowania autoregresyjnego: kaskada musi wytworzyć transkrypcję element po elemencie,
podczas gdy ścieżka latentna wykonuje pojedyncze przejście w przód przez enkoder oraz adapter.
Pomiary te są dla wyboru architektury równorzędne z jakością, ponieważ w zastosowaniu klinicznym o
przydatności systemu decyduje nie tylko trafność wyników, lecz także czas oczekiwania na nie.

\todo{Uzupełnić zestawienie pomiarów zasobów i opóźnień po wykonaniu przebiegów}

\subsection{Wnioski}
Podrozdział zestawia obie architektury wraz z ich osiągami, wskazując warunki, w których przewaga
jednej z nich jest rozstrzygająca, oraz warunki, w których wybór pozostaje kompromisem między
jakością a kosztem obliczeniowym.

\todo{Uzupełnić wnioski po zebraniu wyników z podrozdziałów \ref{sec:czyste}--\ref{sec:zasoby}}
